\section{Related Work}
\label{sec:related-work}

\textbf{RLVR for LLM reasoning.}
Reinforcement learning is a standard post-training tool for aligning and improving LLMs~\cite{christiano2017deep,bai2022training,ouyang2022training,rafailov2023direct}. Recent reasoning systems increasingly rely on reinforcement learning with verifiable rewards (RLVR), where correctness can be checked directly for tasks such as mathematics and code~\cite{guo2025deepseek,hu2025reinforce++,liu2503understanding,jaech2024openai,Zhang2025SRPOAC}. Algorithmic work improves this training loop through value-enhanced PPO variants such as VC-PPO~\cite{yuan2025s} and VAPO~\cite{yue2025vapo}, group-based methods such as RLOO~\cite{ahmadian2024back} and GRPO~\cite{shao2024deepseekmath}, replay-based updates~\cite{liang2025squeeze,li2025repo,zhang2025rlep}, and objective modifications such as Dr. GRPO~\cite{liu2503understanding} and GSPO~\cite{zheng2025group}. These methods improve optimization stability or sample reuse, but they do not directly solve the pre-rollout allocation problem studied here: which prompts deserve expensive response generation at the current training step.

\textbf{Prompt selection and data-efficient RL training.}
Data curation is essential for efficient LLM training. Offline filtering methods rank prompts before training using quality, difficulty, diversity, or scale heuristics~\cite{Wang2025ReinforcementLF,Fatemi2025ConciseRV,Li2025LIMRLI,shi2025efficient,tang2025towards,zhou2023lima,ye2025limo}. They are efficient once training begins, but they are policy-agnostic and cannot adapt as the model's competence changes. Rollout-based online methods, including Dynamic Sampling in DAPO~\cite{yu2025dapo} and related online down-sampling or difficulty-filtering approaches~\cite{meng2025mm,foster2025learning,xu2025not,lin2025cppo,sun2025efficient,bae2025online}, are more precise because they observe current rewards, but they require the rollouts they aim to save. History-based and estimator-based approaches such as GRESO~\cite{zheng2025act}, Bayesian task selection~\cite{chen2025self,zeng2025cures,shen2025bots}, and curriculum estimators use past trajectories or uncertainty models to reduce selection cost, but their estimates can become stale in a changing policy environment. Auxiliary-model methods such as DOTS~\cite{sun2025dots} and PCL~\cite{Gao2025PCL} add external or on-policy predictors for difficulty, improving adaptivity at the cost of additional compute or memory.

Overall, prior methods face limitations in online and efficient prompt selection. HIVE differs by explicitly decomposing the problem into history-informed coarse narrowing and prompt-side online verification, allowing it to retain the efficiency of historical filtering while correcting stale metadata before rollout. Detailed related work is provided in Appendix~\ref{appendix-sec:detailed_related_work}.
